\clearpage\thispagestyle{empty}
\pdfbookmark[0]{Quellen zum Juno-Bericht vom Dezember 1804}{juno180412-sources}
\begin{center}\textbf{Quellen zum Juno-Bericht vom Dezember 1804}\end{center}
\noindent Diese KI-Übertragung in heutiges Deutsch folgt dem Auszug aus Franz Xaver von Zachs Bericht in \href{https://gdz.sub.uni-goettingen.de/id/PPN236007467}{Gauß’ \textit{Werken}, Band~VI} (Göttingen, 1874), S.~258--260. Die entsprechende Passage in \textit{Ueber die Juno}, Beitrag LII der \textit{Monatlichen Correspondenz}, Band~X, Dezember 1804, beginnt auf \href{https://www.digitale-sammlungen.de/en/view/bsb10538602?page=566}{S.~552} und reicht bis zu den Seeberger Beobachtungen auf \href{https://www.digitale-sammlungen.de/en/view/bsb10538602?page=569}{S.~555}.
\par\smallskip\noindent Von Zach zitiert Carl Friedrich Gauß zu den dritten Elementen. Die Tabellen enthalten unter anderem Beobachtungen, die von Zach auf dem Seeberg und Nevil Maskelyne in Greenwich anstellten. Die Zeitschrift enthält die Vergleichstabellen auf \href{https://www.digitale-sammlungen.de/en/view/bsb10538602?page=567}{S.~553} und die Ephemeride auf \href{https://www.digitale-sammlungen.de/en/view/bsb10538602?page=568}{S.~554}. Die Werkausgabe beschreibt die Ephemeride in eckigen Klammern.
\par\smallskip\noindent Übertragung und Satz wurden mit KI erstellt. Die Digitalisate der Zeitschrift stammen von der Bayerischen Staatsbibliothek.
\clearpage
